\documentclass[aspectratio=169]{beamer} 
\synctex=1
% 16:9 宽屏比例

\usepackage{../../Styles/mystyle-PPT}

\graphicspath{{./image/}}% 指定图片目录，后续可以直接使用图片文件名。

% beamer 主题、颜色与字体
\usetheme{metropolis}
\definecolor{mycolor}{RGB}{0, 176, 183}% 自定义青色
\definecolor{mycolor2}{RGB}{0, 191, 165}% 自定义亮绿色
% \usecolortheme[named=mycolor]{structure}% 让所有结构性元素（例如\item的序号等）都默认保持mycolor颜色

\usefonttheme{professionalfonts}% 让所有公式的字母、数字恢复为标准 LaTeX（Computer Modern）的衬线数学字体

\makeatother

\begin{document}

% 标题页
\begin{frame}[plain]
\vfill
\begin{center}
{\Huge\bfseries\color{mycolor} 2016 年全国硕士研究生招生考试}\\[0.5em]
{\Huge\bfseries\color{mycolor} 数学一试题解析}\\[2em]
\rule{0.8\textwidth}{2pt}\\[2em]
{\large 数学二班第六组}\\[0.6em]
{\normalsize 组长：邹文杰\quad 组员：褚涵、刘川蓉、王美婷、祖文颜、付佳一}\\[1.5em]
{\normalsize\color{gray}\today}
\end{center}
\vfill
\end{frame}

% 目录
\begin{frame}{目录 CONTENTS}
\begin{columns}[T]
\column{0.48\textwidth}
\begin{block}{选择题}
1--8 小题
\end{block}
\begin{block}{填空题}
9--14 小题
\end{block}
\column{0.48\textwidth}
\begin{block}{解答题}
15--23 小题
\end{block}
\end{columns}
\end{frame}

% ============================================================
% 选择题部分
% ============================================================
\section{选择题}

\begin{frame}{第1题}
\begin{block}{题目}
若反常积分 \(\displaystyle\int_{0}^{+\infty} \frac{1}{x^{a}(1+x)^{b}} \,\mathrm{d}x\) 收敛，则
$\mathrm{A.}\ a<1 \text{ 且 } b>1; \quad
\mathrm{B.}\ a>1 \text{ 且 } b>1; \quad
\mathrm{C.}\ a<1 \text{ 且 } a+b>1; \quad
\mathrm{D.}\ a>1 \text{ 且 } a+b>1$.
\\
\textcolor{red}{\textbf{答案：}$\mathbf{C}$} \\
\textbf{解析：}分段判断收敛性，积分分成$0$到$1$和$1$到$+\infty$两段.$x\to0^+$ 时要求$a<1$，$x\to+\infty$ 时要求 $a+b>1$.
\end{block}
\end{frame}

\begin{frame}{第2题}
\begin{block}{题目}
已知函数 \(f(x)= \begin{cases}2(x-1), & x<1, \\ \ln x, & x \geqslant 1,\end{cases}\)
则 \(f(x)\) 的一个原函数是
\begin{enumerate}[A.]
\item \(F(x)=\begin{cases}(x-1)^{2}, & x<1, \\ x(\ln x-1), & x \geqslant 1;\end{cases}\)
\item \(F(x)=\begin{cases}(x-1)^{2}, & x<1, \\ x(\ln x+1)-1, & x \geqslant 1;\end{cases}\)
\item \(F(x)=\begin{cases}(x-1)^{2}, & x<1, \\ x(\ln x+1)+1, & x \geqslant 1;\end{cases}\)
\item \(F(x)=\begin{cases}(x-1)^{2}, & x<1, \\ x(\ln x-1)+1, & x \geqslant 1.\end{cases}\)
\end{enumerate}
\textcolor{red}{\textbf{答案：}$\mathbf{D}$} \\
\textbf{解析：}原函数在 $x=1$ 处连续，代入验证左右值相等即可.
\end{block}
\end{frame}

\begin{frame}{第3题}
\begin{block}{题目}
若 \(y_1=(1+x^{2})^{2}-\sqrt{1+x^{2}}\)，\(y_2=(1+x^{2})^{2}+\sqrt{1+x^{2}}\) 是微分方程 \(y'+p(x) y=q(x)\) 的两个解，则 \(q(x)=\)
$\mathrm{A.}\ 3 x\left(1+x^{2}\right); \quad
\mathrm{B.}\ -3 x\left(1+x^{2}\right); \quad
\mathrm{C.}\ \dfrac{x}{1+x^{2}}; \quad
\mathrm{D.}\ -\dfrac{x}{1+x^{2}}$.
\\
\textcolor{red}{\textbf{答案：}$\mathbf{A}$} \\
\textbf{解析：}两解之差$y_2-y_1=2\sqrt{1+x^2}$为齐次解，代入齐次方程得
\begin{align*}
\left( 2\sqrt{1+x^2} \right) ^{\prime}+2p\left( x \right) \sqrt{1+x^2}=0\Longrightarrow p\left( x \right) =-\frac{x}{1+x^2}.
\end{align*}
再代回原方程得
\begin{align*}
q\left( x \right) &=y_{2}^{\prime}-\frac{x}{1+x^2}y_2=4x\left( 1+x^2 \right) +\frac{x}{\sqrt{1+x^2}}-\frac{x}{1+x^2}\left[ \left( 1+x^2 \right) ^2+\sqrt{1+x^2} \right] 
\\
&=3x\left( 1+x^2 \right) .
\end{align*}
\end{block}
\end{frame}

\begin{frame}{第4题}
\begin{block}{题目}
已知函数
\[f(x)=\begin{cases} x, & x \leqslant 0, \\[4pt]
\dfrac{1}{n}, & \dfrac{1}{n+1}<x \leqslant \dfrac{1}{n},\ n=1,2,\cdots,\end{cases}\]
则
\begin{enumerate}[A.]
\item \(x=0\) 是 \(f(x)\) 的第一类间断点;
\item \(x=0\) 是 \(f(x)\) 的第二类间断点;
\item \(f(x)\) 在 \(x=0\) 处连续但不可导;
\item \(f(x)\) 在 \(x=0\) 处可导.
\end{enumerate}
\end{block}
\end{frame}

\begin{frame}{第4题解答}
\begin{block}{}
\textcolor{red}{\textbf{答案：}$\mathbf{D}$} \\
\textbf{解析：}左右极限显然均为0，连续；左导数显然为1.考虑右导数.
对$\forall x\in \left( 0,1 \right) ,$存在$n_x\in \mathbb{N} ,$使得$\frac{1}{n_x+1}<x\leqslant \frac{1}{n_x}.$从而
\begin{align*}
\frac{f\left( x \right)}{x}=\frac{1}{n_xx}<\frac{1}{n_x\left( n_x+1 \right)}.
\end{align*}
令$x\rightarrow 0^+,$此时$n_x\rightarrow +\infty .$故
\begin{align*}
\varlimsup\limits_{x\rightarrow 0^+}\frac{f\left( x \right)}{x}\leqslant \varlimsup\limits_{n_x\rightarrow +\infty}\frac{1}{n_x\left( n_x+1 \right)}=0\Longrightarrow \lim\limits_{x\rightarrow 0^+}\frac{f\left( x \right)}{x}=0.
\end{align*}
\end{block}
\end{frame}

\begin{frame}{第5题}
\begin{block}{题目}
设矩阵 \(A\) 是可逆矩阵，且 \(A\) 与 \(B\) 相似，则下列结论错误的是
\begin{enumerate}[A.]
\item \(A^{T}\) 与 \(B^{T}\) 相似;
\item \(A^{-1}\) 与 \(B^{-1}\) 相似;
\item \(A+A^{T}\) 与 \(B+B^{T}\) 相似;
\item \(A+A^{-1}\) 与 \(B+B^{-1}\) 相似.
\end{enumerate}
\textcolor{red}{\textbf{答案：}$\mathbf{C}$} \\
\textbf{解析：}相似变换下转置的过渡矩阵对应 $P^T$ 而非 $P^{-1}$，故 $A+A^T$ 未必与 $B+B^T$ 相似.
\end{block}
\end{frame}

\begin{frame}{第5题反例}
反例：取
\[
A=\begin{pmatrix}1&1&0\\0&1&0\\0&0&2\end{pmatrix},\quad
B=\begin{pmatrix}1&0&0\\0&1&1\\0&0&2\end{pmatrix},
\]
二者相似（均为若尔当块 \(J_2(1)\oplus J_1(2)\)），但
\[
A+A^T=\begin{pmatrix}2&1&0\\1&2&0\\0&0&4\end{pmatrix},\quad
B+B^T=\begin{pmatrix}2&0&0\\0&2&1\\0&1&4\end{pmatrix},
\]
前者特征值为 \(2,2,4\)，后者特征值为 \(2,3\pm\sqrt{2}\)，特征值不同，故不相似。
\end{frame}

\begin{frame}{第6题}
\begin{block}{题目}
设二次型 \(f(x_{1}, x_{2}, x_{3})=x_{1}^{2}+x_{2}^{2}+x_{3}^{2}+4 x_{1} x_{2}+4 x_{1} x_{3}+4 x_{2} x_{3}\)，则 \(f(x_{1}, x_{2}, x_{3})=2\) 在空间直角坐标下表示的二次曲面为
\begin{enumerate}[A.]
\item 单叶双曲面;
\item 双叶双曲面;
\item 椭球面;
\item 柱面.
\end{enumerate}
\end{block}
\end{frame}

\begin{frame}{第6题解答}
\begin{block}{}
\textcolor{red}{\textbf{答案：}$\mathbf{B}$} \\
\textbf{解析：}$f\left( \boldsymbol{x} \right) =\boldsymbol{x}^TA\boldsymbol{x}$,其中$\boldsymbol{x}=(x_1,x_2,x_3)^T,A=\left( \begin{matrix}
1&		2&		2\\
2&		1&		2\\
2&		2&		1\\
\end{matrix} \right) $.从而
\begin{align*}
\left| \lambda I-A \right|&=\left| \begin{matrix}
\lambda -1&		-2&		-2\\
-2&		\lambda -1&		-2\\
-2&		-2&		\lambda -1\\
\end{matrix} \right|=\left| \begin{matrix}
\lambda -1&		-2&		-2\\
-\lambda -1&		\lambda +1&		0\\
-\lambda -1&		0&		\lambda +1\\
\end{matrix} \right|
\\
&=\left( \lambda -1 \right) \left( \lambda +1 \right) ^2-\left( \lambda +1 \right) \left( -2 \right) \left( -\lambda -1 \right) -\left( \lambda +1 \right) \left( -2 \right) \left( -\lambda -1 \right) 
\\
&=\left( \lambda -1 \right) \left( \lambda +1 \right) ^2-4\left( \lambda +1 \right) ^2=\left( \lambda -5 \right) ^2\left( \lambda +1 \right) ^2.
\end{align*}
因此$A$的特征值为$5,-1,-1$，故$f\left( \boldsymbol{x} \right)$的标准型为 $5y_1^2-y_2^2-y_3^2=2$，对应双叶双曲面.
\end{block}
\end{frame}

\begin{frame}{第7题}
\begin{block}{题目}
设随机变量 \(X \sim N(\mu, \sigma^{2})\ (\sigma>0)\)，记 \(p=P\{X \leqslant \mu+\sigma^{2}\}\)，则
\begin{enumerate}[A.]
\item \(p\) 随着 \(\mu\) 的增加而增加;
\item \(p\) 随着 \(\sigma\) 的增加而增加;
\item \(p\) 随着 \(\mu\) 的增加而减少;
\item \(p\) 随着 \(\sigma\) 的增加而减少.
\end{enumerate}
\textcolor{red}{\textbf{答案：}$\mathbf{B}$} \\
\textbf{解析：}令$y=\frac{X-\mu}{\sigma}$,则$Y\sim N(0,1),p=P\{Y\leqslant \sigma\}$.
显然$p$ 随 $\sigma$ 增大而增大.
\end{block}
\end{frame}

\begin{frame}{第8题}
\begin{block}{题目}
随机试验 \(E\) 有三种两两不相容的结果 \(A_1,A_2,A_3\)，且发生概率均为 \(\dfrac13\)。将试验独立重复2次，\(X\) 表示 \(A_1\) 发生的次数，\(Y\) 表示 \(A_2\) 发生的次数，则 \(\rho_{XY}=\) \\
$\mathrm{A.}\ -\dfrac{1}{2};\quad \mathrm{B.}\ -\dfrac{1}{3};\quad \mathrm{C.}\ \dfrac{1}{3};\quad \mathrm{D.}\ \dfrac{1}{2}.$

\textcolor{red}{\textbf{答案：} \(\mathbf{A}\)}\\
\textbf{解析：}
列出\((X,Y)\)的所有可能结果和概率：\((2,0),(0,2)\) 概率各 \(\frac{1}{9}\)；\((1,1)\) 概率 \(\frac{2}{9}\)；\((1,0),(0,1)\) 概率各 \(\frac{2}{9}\)；\((0,0)\) 概率 \(\frac{1}{9}\).

期望：\(E(X)=E(Y)=\frac23\)，\(E(X^2)=E(Y^2)=\frac89\),且 \(E(XY)=\frac29\).

协方差：\(\operatorname{Cov}(X,Y)=E(XY)-E(X)E(Y)=\frac29-\frac23\cdot\frac23=-\frac29\).

方差：\(D(X)=E(X^2)-[E(X)]^2=\frac49,
D(Y)=E(Y^2)-[E(Y)]^2=\frac49\).

故 \(\rho _{XY}=\frac{\mathrm{Cov}\left( X,Y \right)}{\sqrt{D\left( X \right) D\left( Y \right)}}=-\frac{1}{2}.\)
\end{block}
\end{frame}

% ============================================================
% 填空题部分
% ============================================================
\section{填空题}

\begin{frame}{第9题}
\begin{block}{题目}
\(\displaystyle\lim _{x \to 0} \frac{\int_{0}^{x} t \ln (1+t \sin t) \,\mathrm{d} t}{1-\cos x^{2}}=\) \underline{$\quad\textcolor{red}{\dfrac{1}{2}}\quad$} 
\end{block}

\begin{block}{解析}
\begin{align*}
&\lim_{x\rightarrow 0} \frac{\int_0^x{t\ln \left( 1+t\sin t \right) \mathrm{d}t}}{1-\cos x^2}\xlongequal{\mathrm{L}' \mathrm{Hospital}}\lim_{x\rightarrow 0} \frac{x\ln \left( 1+x\sin x \right)}{2x\sin x^2}\\
&=\lim_{x\rightarrow 0} \frac{x^2\sin x}{2x^3}
=\lim_{x\rightarrow 0} \frac{\sin x}{2x}=\frac{1}{2}.
\end{align*}
\end{block}
\end{frame}

\begin{frame}{第10题}
\begin{block}{题目}
向量场 \(\boldsymbol{A}(x, y, z)=(x+y+z)\boldsymbol{i}+xy\boldsymbol{j}+z\boldsymbol{k}\) 的旋度 \(\operatorname{rot}\boldsymbol{A}=\) \underline{$\quad\textcolor{red}{\boldsymbol{j}+(y-1)\boldsymbol{k}}\quad$} 
\end{block}

\begin{block}{解析}
由旋度公式 $\operatorname{rot}\boldsymbol{A}=\nabla\times \boldsymbol{A}=\begin{vmatrix}\boldsymbol{i}&\boldsymbol{j}&\boldsymbol{k}\\[4pt]\dfrac{\partial}{\partial x}&\dfrac{\partial}{\partial y}&\dfrac{\partial}{\partial z}\\[4pt]P&Q&R\end{vmatrix}$,将$P=x+y+z,q=xy,R=z$代入计算得 $\operatorname{rot}\boldsymbol{A}=\boldsymbol{j}+(y-1)\boldsymbol{k}$.
\end{block}
\end{frame}

\begin{frame}{第11题}
\begin{block}{题目}
设函数 \(z=z(x, y)\) 由方程 \((x+1) z-y^{2}=x^{2} f(x-z, y)\) 确定，则 \(\left.\mathrm{d} z\right|_{(0,1)}=\) \underline{$\quad \textcolor{red}{-\mathrm{d} x+2\mathrm{d} y}\quad$} 
\end{block}

\begin{block}{解析}
代入 $(x,y)=(0,1)$ 得 $z(0,1)=1$.方程两边分别对$x,y$求偏导得
\begin{gather*}
\left( x+1 \right) \frac{\partial z}{\partial x}+z=2xf\left( x-z,y \right) +x^2\frac{\partial f}{\partial x}\left( 1-\frac{\partial z}{\partial x} \right) ,
\\
\left( x+1 \right) \frac{\partial z}{\partial y}-2y= -x^2\frac{\partial f}{\partial y}\frac{\partial z}{\partial y}.
\end{gather*}
并代入$(x,y)=(0,1)$得 $\left.\dfrac{\partial z}{\partial x}\right|_{(0,1)}=-1,\left.\dfrac{\partial z}{\partial y}\right|_{(0,1)}=2$，故 $\mathrm{d}z=-\mathrm{d}x+2\mathrm{d}y$.
\end{block}
\end{frame}

\begin{frame}{第12题}
\begin{block}{题目}
设函数 \(f(x)=\arctan x-\dfrac{x}{1+a x^{2}}\)，若 \(f^{\prime \prime \prime}(0)=1\)，则 \(a=\) \underline{$\quad \textcolor{red}{\dfrac{1}{2}}\quad$} 
\end{block}

\begin{block}{解析}
利用带Peano余项的Taylor展开公式将$f$在$x=0$处展开得
\begin{align*}
f(x)&=x-\frac{x^3}{3}+o(x^3)+x-ax^3+o(x^3)
\\
&=\left( a-\frac{1}{3} \right) x^3+o(x^3),\quad x\rightarrow 0.
\end{align*}
故$f'''(0)=3!\left(a-\dfrac{1}{3}\right)=1$,因此 $a=\dfrac{1}{2}$.
\end{block}
\end{frame}

\begin{frame}{第13题}
\begin{block}{题目}
行列式 \(\begin{vmatrix}
\lambda & -1 & 0 & 0 \\
0 & \lambda & -1 & 0 \\
0 & 0 & \lambda & -1 \\
4 & 3 & 2 & \lambda+1
\end{vmatrix} =\) \underline{$\quad \textcolor{red}{\lambda^{4}+\lambda^{3}+2\lambda^{2}+3\lambda+4} \quad$} 
\end{block}

\begin{block}{解析}
按第一列展开得
\begin{align*}
\left| \begin{matrix}
\lambda&		-1&		0&		0\\
0&		\lambda&		-1&		0\\
0&		0&		\lambda&		-1\\
4&		3&		2&		\lambda +1\\
\end{matrix} \right|&=\lambda \left| \begin{matrix}
\lambda&		-1&		0\\
0&		\lambda&		-1\\
3&		2&		\lambda +1\\
\end{matrix} \right|-4\left| \begin{matrix}
-1&		0&		0\\
\lambda&		-1&		0\\
0&		\lambda&		-1\\
\end{matrix} \right|\\
&=\lambda ^2\left| \begin{matrix}
\lambda&		-1\\
2&		\lambda +1\\
\end{matrix} \right|+3\lambda \left| \begin{matrix}
-1&		0\\
\lambda&		-1\\
\end{matrix} \right|+4
\\
&=\lambda ^3\left( \lambda +1 \right) +2\lambda ^2+3\lambda +4=\lambda ^4+\lambda ^3+2\lambda ^2+3\lambda +4.
\end{align*}
\end{block}
\end{frame}

\begin{frame}{第14题}
\begin{block}{题目}
设 \(X_{1}, X_{2}, \cdots, X_{n}\) 为来自总体 \(N(\mu, \sigma^{2})\) 的简单随机样本，样本均值 \(\bar{X}=9.5\)，$\mu$ 的置信度0.95双侧置信上限为10.8，则置信区间为 \underline{$\quad \textcolor{red}{(8.2,\ 10.8)}\quad$} 
\end{block}

\begin{block}{解析}
正态总体双侧置信区间关于样本均值对称，下限为 $9.5-(10.8-9.5)=8.2$，故置信区间为 $(8.2,\,10.8)$.
\end{block}
\end{frame}

% ============================================================
% 解答题部分
% ============================================================
\section{解答题}

\begin{frame}{第15题}
\begin{block}{题目}
已知平面区域 \(D=\{(r, \theta) \mid 2 \leqslant r \leqslant 2(1+\cos \theta),\ -\dfrac{\pi}{2} \leqslant \theta \leqslant \dfrac{\pi}{2}\}\)，计算二重积分 \(\displaystyle\iint_{D} x\,\mathrm{d}\sigma\).
\end{block}
\begin{block}{解答}
极坐标下 \(x=r\cos\theta,\ \mathrm{d}\sigma=r\mathrm{d}r\mathrm{d}\theta\)，利用对称性
\begin{align*}
\iint_D x\,\mathrm{d}\sigma &= \int_{-\frac{\pi}{2}}^{\frac{\pi}{2}} \cos\theta \,\mathrm{d}\theta \int_{2}^{2(1+\cos\theta)} r^2 \,\mathrm{d}r \\
&= \frac{16}{3}\int_{0}^{\frac{\pi}{2}} \cos\theta\big[(1+\cos\theta)^3-1\big]\mathrm{d}\theta \\
&= \frac{16}{3}\int_{0}^{\frac{\pi}{2}} \big(3\cos^2\theta+3\cos^3\theta+\cos^4\theta\big)\mathrm{d}\theta \\
&= 5\pi+\frac{32}{3}.
\end{align*}
\end{block}
\end{frame}

\begin{frame}{第16题}
\begin{block}{题目}
设函数 \(y(x)\) 满足方程 \(y''+2 y'+k y=0\)，其中 \(0<k<1\).
\begin{enumerate}[(1)]
\item 证明：反常积分 \(\displaystyle\int_{0}^{+\infty} y(x) \,\mathrm{d}x\) 收敛；

\item 若 \(y(0)=1,\ y'(0)=1\)，求 \(\displaystyle\int_{0}^{+\infty} y(x) \,\mathrm{d}x\) 的值.
\end{enumerate} 
\end{block}
\end{frame}

\begin{frame}
\begin{block}{解答}
\begin{enumerate}[(1)]
\item 特征方程 \(r^2+2r+k=0\)，根为 \(r_1=-1+\sqrt{1-k},r_2=-1-\sqrt{1-k}\)，均为负实根，故 \(y=C_1e^{r_1x}+C_2e^{r_2x}\),其中$C_1,C_2$为任意常数,因此积分收敛.

\item 方程两边在 \([0,+\infty)\) 积分
\begin{align*}
\int_0^{+\infty} y''\mathrm{d}x + 2\int_0^{+\infty} y'\mathrm{d}x + k\int_0^{+\infty} y\mathrm{d}x = 0,
\end{align*}
代入边界值得 \(-1-2 + kI=0\)，故 \(I=\dfrac{3}{k}\).
\end{enumerate}    
\end{block}
\end{frame}

\begin{frame}{第17题}
\begin{block}{题目}
设函数 \(f(x, y)\) 满足 \(\dfrac{\partial f(x, y)}{\partial x}=(2 x+1) e^{2 x-y}\)，且 \(f(0, y)=y+1\)，\(L_t\) 是从 \((0,0)\) 到 \((t,0)\) 再到 \((t,1)\) 的折线段. 计算曲线积分 \(I(t)=\displaystyle\int_{L_{t}} \frac{\partial f}{\partial x} \mathrm{d} x+\frac{\partial f}{\partial y} \mathrm{d} y\)，并求 \(I(t)\) 的最小值.
\end{block}
\begin{block}{解答}
对 $x$ 积分并代入初值得 \(f(x,y)=x e^{2x-y}+y+1\)，积分与路径无关，故
\begin{align*}
I(t)=\int_{L_t}{\mathrm{d}f}=f(t,1)-f(0,0)=t e^{2t-1}+1.
\end{align*}
求导得 \(I'(t)=(2t+1)e^{2t-1}\)，令 \(I'(t)=0\) 得 \(t=-\dfrac12\)，
最小值为 \(I\!\left(-\dfrac12\right)=1-\dfrac{1}{2e^2}\).
\end{block}
\end{frame}

\begin{frame}{第18题}
\begin{block}{题目}
设有界区域 \(\Omega\) 由平面 \(2 x+y+2 z=2\) 与三个坐标平面围成，\(\Sigma\) 为 \(\Omega\) 整个表面的外侧，计算曲面积分
\[I=\iint_{\Sigma}\left(x^{2}+1\right) \mathrm{d} y \mathrm{d} z-2 y \,\mathrm{d} z \mathrm{d} x+3 z \,\mathrm{d} x \mathrm{d} y.\]
\end{block}
\begin{block}{解答}
由Gauss公式及\(\operatorname{div}\boldsymbol{F}=\nabla\cdot \boldsymbol{F}=2x+1\)，故
\begin{align*}
&I=\iiint_{\Omega}{\left( 2x+1 \right) \,\mathrm{d}V}=\int_0^1{\mathrm{d}x\int_0^{2-2x}{\mathrm{d}y\int_0^{1-x-\frac{y}{2}}{\left( 2x+1 \right) \mathrm{d}z}}}
\\
&=\int_0^1{\mathrm{d}x\int_0^{2-2x}{\left( 2x+1 \right) \left( 1-x-\frac{y}{2} \right) \mathrm{d}y}}
=\int_0^1{\left( 2x+1 \right) \left( 1-x \right) ^2\mathrm{d}x}=\frac{1}{2}.
\end{align*}
\end{block}
\end{frame}

\begin{frame}{第19题}
\begin{block}{题目}
已知函数 \(f(x)\) 可导，且 \(f(0)=1,\ 0<f'(x)<\dfrac{1}{2}\). 设数列 \(\{x_{n}\}\) 满足 \(x_{n+1}=f(x_{n})\ (n=1,2,\cdots)\). 证明：
(1) 级数 \(\displaystyle\sum_{n=1}^{\infty}(x_{n+1}-x_{n})\) 绝对收敛；
(2) \(\displaystyle\lim_{n \to \infty} x_{n}\) 存在，且 \(0<\lim_{n \to \infty} x_{n}<2\).
\end{block}
\begin{block}{证明}
(1) 由Lagrange中值定理
\[|x_{n+1}-x_n|=|f(x_n)-f(x_{n-1})|=|f'(\xi)|\cdot|x_n-x_{n-1}|<\frac12|x_n-x_{n-1}|,\]
由比值判别法知级数绝对收敛.

(2) 级数绝对收敛故 \(\{x_n\}\) 收敛，设极限为 $a$，则 $a=f(a)$.
令 $g(x)=f(x)-x$，$g(0)=1>0,\ g(2)<0$，由零点定理与单调性得 $0<a<2$.
\end{block}
\end{frame}

\begin{frame}{第20题}
\begin{block}{题目}
设矩阵 \(A=\begin{pmatrix}1&-1&-1\\2&a&1\\-1&1&a\end{pmatrix}\)，\(B=\begin{pmatrix}2&2\\1&a\\-a-1&-2\end{pmatrix}\). 当 $a$ 为何值时，方程 \(AX=B\) 无解、有唯一解、有无穷多解？在有解时求解.
\end{block}

\begin{block}{解答：增广矩阵的初等行变换}
对增广矩阵 \((A|B)\) 做初等行变换：
\[
(A|B) = \begin{pmatrix}
1 & -1 & -1 & \vert & 2 & 2 \\
2 & a & 1 & \vert & 1 & a \\
-1 & 1 & a & \vert & -a-1 & -2
\end{pmatrix}\xrightarrow[R_3+R_1]{R_2-2R_1} \begin{pmatrix}
1 & -1 & -1 & \vert & 2 & 2 \\
0 & a+2 & 3 & \vert & -3 & a-4 \\
0 & 0 & a-1 & \vert & 1-a & 0
\end{pmatrix}.
\]
\end{block}
\end{frame}

\begin{frame}{第20题：分类讨论（1）无解情况}
\begin{block}{情况一：当 \(a=-2\) 时}
将 \(a=-2\) 代入变换后的矩阵：
\[
(A|B) = \begin{pmatrix}
1 & -1 & -1 & \vert & 2 & 2 \\
0 & 0 & 3 & \vert & -3 & -6 \\
0 & 0 & -3 & \vert & 3 & 0
\end{pmatrix}
\xrightarrow{R_3+R_2} \begin{pmatrix}
1 & -1 & -1 & \vert & 2 & 2 \\
0 & 0 & 3 & \vert & -3 & -6 \\
0 & 0 & 0 & \vert & 0 & -6
\end{pmatrix}
\]
此时，\(r(A) = 2\)，而 \(r(A|B) = 3\)，两者不相等。

\textbf{结论：} 当 \(a=-2\) 时，方程 \(AX=B\) \textbf{无解}。
\end{block}
\end{frame}

\begin{frame}{第20题：分类讨论（2）唯一解情况}
\begin{block}{情况二：当 \(a\neq-2\) 且 \(a\neq1\) 时}
此时 \(a+2 \neq 0\) 且 \(a-1 \neq 0\)， \(r(A)=3\)且 \(r(A|B)=3\)，方程有唯一解。
继续对原行变换后的矩阵进行求解得:$x_1=1,x_2=0,x_3=-1$和$x_1=\frac{3a}{a+2},x_2=\frac{a-4}{a+2},x_3=0$分别对应$B$的第一列和第二列.

\textbf{结论：} 当 \(a\neq-2\) 且 \(a\neq1\) 时，方程有唯一解：
\[
X = \begin{pmatrix}1 & \dfrac{3a}{a+2} \\[6pt] 0 & \dfrac{a-4}{a+2} \\[6pt] -1 & 0\end{pmatrix}
\]
\end{block}
\end{frame}

\begin{frame}{第20题：分类讨论（3）无穷多解情况}
\begin{block}{情况三：当 \(a=1\) 时}
将 \(a=1\) 代入变换后的矩阵：
\[
(A|B) = \begin{pmatrix}
1 & -1 & -1 & \vert & 2 & 2 \\
0 & 3 & 3 & \vert & -3 & -3 \\
0 & 0 & 0 & \vert & 0 & 0
\end{pmatrix}
\]
此时，\(r(A) = r(A|B) = 2 < 3\)，方程组有无穷多解。
继续进行初等行变换化简化行阶梯形：
\[
(A|B) \xrightarrow{R_2 \times \frac{1}{3}} \begin{pmatrix}
1 & -1 & -1 & \vert & 2 & 2 \\
0 & 1 & 1 & \vert & -1 & -1 \\
0 & 0 & 0 & \vert & 0 & 0
\end{pmatrix}
\xrightarrow{R_1 + R_2} \begin{pmatrix}
1 & 0 & 0 & \vert & 1 & 1 \\
0 & 1 & 1 & \vert & -1 & -1 \\
0 & 0 & 0 & \vert & 0 & 0
\end{pmatrix}
\]
\end{block}
\end{frame}

\begin{frame}{第20题：分类讨论（3）无穷多解情况}
对应的方程组为：\(\begin{cases} x_1 = 1 \\ x_2 + x_3 = -1 \end{cases}\)
取自由变量 \(x_3 = c\)，则 \(x_2 = -1 - c\)。
对于第一列，令 \(x_3 = c_1\)；对于第二列，令 \(x_3 = c_2\)。

\textbf{结论：} 当 \(a=1\) 时，方程有无穷多解，通解为：
\[
X=\begin{pmatrix}1 & 1\\ -1 & -1\\ 0 & 0\end{pmatrix}+\begin{pmatrix}0 & 0\\ -c_1 & -c_2\\ c_1 & c_2\end{pmatrix},\quad c_1,c_2\in\mathbb{R}
\]
\end{frame}
\begin{frame}{第21题}
\begin{block}{题目}
已知矩阵 \(A=\begin{pmatrix}0 & -1 & 1 \\ 2 & -3 & 0 \\ 0 & 0 & 0\end{pmatrix}\).
\begin{enumerate}[(1)]
\item 求 \(A^{99}\)；
\item 设3阶矩阵 \(B=(\alpha_1,\alpha_2,\alpha_3)\) 满足 \(B^2=BA\)，记 \(B^{100}=(\beta_1,\beta_2,\beta_3)\)，将 \(\beta_1,\beta_2,\beta_3\) 表示为 \(\alpha_1,\alpha_2,\alpha_3\) 的线性组合.
\end{enumerate}
\end{block}
\end{frame}

\begin{frame}{第21题 (1)：求特征值与特征向量}
\begin{block}{解答}
计算矩阵 \(A\) 的特征多项式：
\[
|\lambda I - A| = \begin{vmatrix} \lambda & 1 & -1 \\ -2 & \lambda+3 & 0 \\ 0 & 0 & \lambda \end{vmatrix} = \lambda [\lambda(\lambda+3) + 2] = \lambda(\lambda^2 + 3\lambda + 2) = \lambda(\lambda+1)(\lambda+2)
\]
求得特征值为：\(\lambda_1 = 0\), \(\lambda_2 = -1\), \(\lambda_3 = -2\)。

分别求对应的特征向量：
当 \(\lambda_1 = 0\) 时，解 \(Ax=0\)：\(\begin{cases} -y+z=0 \\ 2x-3y=0 \end{cases}\)，令 \(y=2\)，得 \(x=3, z=2\)，即 \(p_1 = (3,2,2)^T\)。
\end{block}
\end{frame}

\begin{frame}{第21题 (1)：求特征向量}
\begin{block}{}
当\(\lambda_2 = -1\) 时，解 \((A+I)x=0\)：\(\begin{cases} x-y+z=0 \\ 2x-2y=0 \\ z=0 \end{cases}\)，令 \(y=1\)，得 \(x=1, z=0\)，即 \(p_2 = (1,1,0)^T\)。

当 \(\lambda_3 = -2\) 时，解 \((A+2I)x=0\)：\(\begin{cases} 2x-y+z=0 \\ 2x-y=0 \\ 2z=0 \end{cases}\)，令 \(x=1\)，得 \(y=2, z=0\)，即 \(p_3 = (1,2,0)^T\)。

构造可逆矩阵 \(P\) 和对角矩阵 \(\Lambda\)：$P = \begin{pmatrix} 3 & 1 & 1 \\ 2 & 1 & 2 \\ 2 & 0 & 0 \end{pmatrix}, \quad \Lambda = \begin{pmatrix} 0 & 0 & 0 \\ 0 & -1 & 0 \\ 0 & 0 & -2 \end{pmatrix}.$
\end{block}
\end{frame}

\begin{frame}{第21题 (1)：计算 \(A^{99}\)}
\begin{block}{}
对 \((P|I)\) 做行变换求 \(P^{-1}\)：
\[
(P|I) = \begin{pmatrix} 3 & 1 & 1 & | & 1 & 0 & 0 \\ 2 & 1 & 2 & | & 0 & 1 & 0 \\ 2 & 0 & 0 & | & 0 & 0 & 1 \end{pmatrix} \xrightarrow{变换} (I|P^{-1})
\]
计算得到：
$P^{-1} = \begin{pmatrix} 0 & 0 & 1/2 \\ 2 & -1 & -2 \\ -1 & 1 & 1/2 \end{pmatrix}$
由于 \(A = P \Lambda P^{-1}\)，所以
\[
A^{99}=P \Lambda^{99} P^{-1} = \begin{pmatrix} 0 & -1 & -2^{99} \\ 0 & -1 & -2^{100} \\ 0 & 0 & 0 \end{pmatrix} \begin{pmatrix} 0 & 0 & 1/2 \\ 2 & -1 & -2 \\ -1 & 1 & 1/2 \end{pmatrix} = \begin{pmatrix} 2^{99}-2 & 1-2^{99} & 2-2^{98} \\ 2^{100}-2 & 1-2^{100} & 2-2^{99} \\ 0 & 0 & 0 \end{pmatrix}
\]
\end{block}
\end{frame}

\begin{frame}{第21题 (2)：推导 \(B^{100}\) 的表达式}
\begin{block}{}
已知 \(B^2 = BA\)。我们要找 \(B^n\) 的规律。
假设 \(B^k = B A^{k-1}\) 对小于等于\(k\)的情形都成立，则再由于 \(B^2 = BA\)可得
\[
B^{k+1} = B A^{k-2} (B^2) = B A^{k-2} (BA) = B A^{k-1}
\]
因此，通过数学归纳法可得，对任意 \(k \ge 1\)，均有 \(B^k = B A^{k-1}\) 成立。

故 \(B^{100} = B A^{99}\)。
代入 \(B=(\alpha_1,\alpha_2,\alpha_3)\) 及 \(A^{99}\)得
\[
B^{100} = (\alpha_1,\alpha_2,\alpha_3) \begin{pmatrix} 2^{99}-2 & 1-2^{99} & 2-2^{98} \\ 2^{100}-2 & 1-2^{100} & 2-2^{99} \\ 0 & 0 & 0 \end{pmatrix}
\]
\end{block}
\end{frame}

\begin{frame}{第21题 (2)：最终结果}
\begin{block}{}
将矩阵乘法展开，得到
\begin{align*}
\beta_1 &= (2^{99}-2)\alpha_1 + (2^{100}-2)\alpha_2 \\
\beta_2 &= (1-2^{99})\alpha_1 + (1-2^{100})\alpha_2 \\
\beta_3 &= (2-2^{98})\alpha_1 + (2-2^{99})\alpha_2
\end{align*}
\end{block}
\end{frame}

\begin{frame}{第22题}
\begin{block}{题目}
设二维随机变量 \((X, Y)\) 在区域 \(D=\{(x, y) \mid 0<x<1,\ x^{2}<y<\sqrt{x}\}\) 上服从均匀分布，令
\[U=\begin{cases} 1, & X \leqslant Y, \\ 0, & X>Y. \end{cases}\]
\begin{enumerate}[(1)]
\item 写出 \((X, Y)\) 的概率密度；
\item 问 \(U\) 与 \(X\) 是否相互独立，并说明理由；
\item 求 \(Z=U+X\) 的分布函数 \(F(z)\).
\end{enumerate}
\end{block}
\end{frame}

\begin{frame}{第22题 (1)：求概率密度}
\begin{block}{}
区域 \(D\) 的面积为：
\[
S = \iint_D \mathrm{d}x\mathrm{d}y = \int_{0}^{1} \mathrm{d}x \int_{x^2}^{\sqrt{x}} \mathrm{d}y = \int_{0}^{1} (\sqrt{x} - x^2) \mathrm{d}x = \left[ \frac{2}{3} x^{3/2} - \frac{1}{3} x^3 \right]_{0}^{1} = \frac{2}{3} - \frac{1}{3} = \frac{1}{3}
\]
因为均匀分布，所以其概率密度为常数 \(\frac{1}{S}\)得
\[
f(x,y) = \begin{cases} 3, & (x,y) \in D, \\ 0, & \text{其他}. \end{cases}
\]
\end{block}
\end{frame}

\begin{frame}{第22题 (2)：判断独立性}
\begin{block}{解答}
判断 \(U\) 与 \(X\) 是否独立，需验证是否满足 \(P(U=0, X \le \frac{1}{2}) = P(U=0)P(X \le \frac{1}{2})\)。
注意到\(U=0\) 即 \(X>Y\)时,该区域为 \(x^2 < y < x\),故分别计算联合概率和边缘概率得
\begin{align*}
&P(U=0,X\le \frac{1}{2})=\int_0^{\frac{1}{2}}{\mathrm{d}x\int_{x^2}^x{3\mathrm{d}y}}=3\int_0^{\frac{1}{2}}{(x}-x^2)\mathrm{d}x
\\
&\quad \quad \quad \quad =3\left[ \frac{x^2}{2}-\frac{x^3}{3} \right] _{0}^{\frac{1}{2}}=3\left( \frac{1}{8}-\frac{1}{24} \right) =\frac{1}{4};
\\
&P(U=0)=\int_0^1{\mathrm{d}x\int_{x^2}^x{3\mathrm{d}y}}=3\int_0^1{(x}-x^2)\mathrm{d}x=3\left[ \frac{x^2}{2}-\frac{x^3}{3} \right] _{0}^{1}=\frac{1}{2},
\\
&P(X\le \frac{1}{2})=\int_0^{\frac{1}{2}}{\mathrm{d}x\int_{x^2}^{\sqrt{x}}{3\mathrm{d}y}}=3\int_0^{\frac{1}{2}}{(\sqrt{x}}-x^2)\mathrm{d}x=\frac{1}{\sqrt{2}}-\frac{1}{8}.
\end{align*}

两者乘积：\(\frac{1}{2} \times (\frac{1}{\sqrt{2}} - \frac{1}{8}) = \frac{1}{2\sqrt{2}} - \frac{1}{16}\)。
明显有 \(\frac{1}{4} \neq \frac{1}{2\sqrt{2}} - \frac{1}{16}\)。

\textbf{结论：} \(U\) 与 \(X\) \textbf{不独立}.
\end{block}
\end{frame}

\begin{frame}{第22题 (3)：求分布函数（第一部分）}
\begin{block}{}
由于 \(X \in (0,1)\)，\(U \in \{0,1\}\)，故 \(Z \in (0,2)\)。
分布函数 \(F(z) = P(Z \le z)\)。
\begin{itemize}
\item 当 \(z \le 0\) 时，\(F(z) = 0\)；
\item 当 \(z \ge 2\) 时，\(F(z) = 1\)；
\item 当 \(0 < z \leqslant 1\) 时,若 \(U=1\)，则 \(Z=1+X > 1 \ge z\) 不满足条件。因此\(U=0\),从而此时 \(Z=X \le z\)。
故
\[
F(z)=P\left( X\le z,Y<X \right) =\int_0^z{\mathrm{d}x}\int_{x^2}^x{3\mathrm{d}x}=3\int_0^z{(x}-x^2)\mathrm{d}x=\frac{3}{2}z^2-z^3.
\]
\end{itemize}
\end{block}
\end{frame}

\begin{frame}{第22题 (3)：求分布函数（第二部分）}
\begin{block}{}
当 \(1 < z \le 2\) 时，
分为两种情况：1) \(U=0\)，此时 \(Z=X \le 1 < z\)，恒成立；2) \(U=1\)，此时 \(Z=1+X \le z\)，需要 \(X \le z-1\)。
因此
\[
F(z) = P(U=0) + P(U=1, X \le z-1)
\]
前面已算得 \(P(U=0) = \frac{1}{2}\)。
而 \(P(U=1, X \le z-1)\) 对应区域为$\left\{ \left( x,y \right) \mid x<y<\sqrt{x},x\le z-1 \right\}$,故
\[
P(U=1,X\le z-1)=\int_0^{z-1}{\mathrm{d}x\int_x^{\sqrt{x}}{3\mathrm{d}y}}=3\int_0^{z-1}{\left( \sqrt{x}-x \right) \mathrm{d}x}=2(z-1)^{\frac{3}{2}}-\frac{3}{2}(z-1)^2.
\]
所以此时有
\begin{align*}
F(z)=\frac{1}{2}+2(z-1)^{\frac{3}{2}}-\frac{3}{2}(z^2-2z+1)=2(z-1)^{\frac{3}{2}}-\frac{3}{2}z^2+3z-1.
\end{align*}
\end{block}
\end{frame}

\begin{frame}{第22题 (3)：分布函数}
综上就有
\[
F(z)= \begin{cases} 0, & z \leqslant 0, \\ \dfrac{3}{2} z^{2}-z^{3}, & 0<z \leqslant 1, \\ 2(z-1)^{\frac{3}{2}}-\dfrac{3}{2} z^{2}+3 z-1, & 1<z \leqslant 2, \\ 1, & z>2. \end{cases}
\]
\end{frame}

\begin{frame}{第23题}
\begin{block}{题目}
设总体 \(X\) 的概率密度为
\[f(x ; \theta)=\begin{cases} \dfrac{3 x^{2}}{\theta^{3}}, & 0<x<\theta, \\ 0, & \text{其他}, \end{cases}\]
其中 \(\theta \in(0,+\infty)\) 为未知参数. \(X_{1},X_{2},X_{3}\) 为来自总体 \(X\) 的简单随机样本，令 \(T=\max \left\{X_{1}, X_{2}, X_{3}\right\}\).
\begin{enumerate}[(1)]
\item 求 \(T\) 的概率密度；
\item 确定 \(a\)，使得 \(aT\) 为 \(\theta\) 的无偏估计.
\end{enumerate}
\end{block}
\end{frame}

\begin{frame}{第23题 (1)：求 \(T\) 的概率密度}
\begin{block}{解答}
首先求总体 \(X\) 的分布函数 \(F_X(t)\)：
当 \(0 < t < \theta\) 时,有
\[
F_X(t) = \int_{-\infty}^{t} f(x) \mathrm{d}x = \int_{0}^{t} \frac{3x^2}{\theta^3} \mathrm{d}x = \left[ \frac{x^3}{\theta^3} \right]_0^t = \frac{t^3}{\theta^3}.
\]
当 \(t \le 0\) 时，\(F_X(t)=0\)；当 \(t \ge \theta\) 时，\(F_X(t)=1\)。

由于 \(X_1, X_2, X_3\) 独立同分布，\(T\) 的分布函数为
\[
F_T(t) = P(T \le t) = P(X_1 \le t, X_2 \le t, X_3 \le t) = [F_X(t)]^3.
\]
所以当 \(0 < t < \theta\) 时，\(F_T(t)=\left( \frac{t^3}{\theta ^3} \right) ^3=\frac{t^9}{\theta ^9}\)。
\end{block}
\end{frame}

\begin{frame}{第23题 (1)：求 \(T\) 的概率密度}
再对 \(t\) 求导，得到 \(T\) 的概率密度：
\[
f_T(t) = F_T'(t) = \begin{cases} \dfrac{9 t^{8}}{\theta^{9}}, & 0<t<\theta, \\ 0, & \text{其他}. \end{cases}
\]
\end{frame}

\begin{frame}{第23题 (2)：确定无偏估计系数 \(a\)}
\begin{block}{}
计算 \(T\) 的数学期望 \(E(T)\)：
\begin{align*}
E(T) &= \int_{-\infty}^{+\infty} t f_T(t) dt = \int_{0}^{\theta} t \cdot \frac{9 t^8}{\theta^9} dt = \frac{9}{\theta^9} \int_{0}^{\theta} t^9 dt
\\
&= \frac{9}{\theta^9} \left[ \frac{t^{10}}{10} \right]_0^\theta = \frac{9}{\theta^9} \cdot \frac{\theta^{10}}{10} = \frac{9}{10}\theta.
\end{align*}
要使 \(aT\) 为 \(\theta\) 的无偏估计，即要求 \(E(aT) = \theta\).此时
\[
E(aT) = a E(T) = a \cdot \frac{9}{10}\theta = \theta,
\]
解得
\[
a = \frac{10}{9}.
\]
\end{block}
\end{frame}

\begin{frame}
\centering
\vfill
\Huge \textcolor{mycolor}{感谢观看} \\
\vfill
\end{frame}

\end{document}